\subsection{The exact cokernel and torsion degree range for the characteristic-two product}

\begin{lemma}[Uniform negative-twist bounds on the constructed surfaces]
\label{lem:OC-surface-uniform-negative-bounds}
\status{Proved; independently audited by the root and reduction-descent agents.}
For the very ample surface $(X,A)$ constructed in
Theorem~\ref{thm:OC-surface-coherent-counterexample} in characteristic $p$, one has
\[
 H^1(X,A^{-r})=0\quad(r\geq p^2-p),\qquad
 H^0(X,\Omega_X^1\otimes A^{-r})=0\quad(r\geq p+1).
\]
In characteristic two these bounds are $r\geq2$ and $r\geq3$, respectively.
\end{lemma}
\begin{proof}
Use the degree-$\ell=p+1$ Raynaud surface $f:S\to C$, its ample generated $L=\OO_S(pE)\otimes f^*N$, and the subsequent multiroot cover $\rho:X\to S$ with $A=\rho^*L$.  Let $P=\mathbf P_C(V)$, where $0\to\OO_C\to V\to T=N^{p+1}\to0$ is the Tango extension.  For $0\leq i\leq p$ define
\[
                  m_i(r)=\left\lceil\frac{p(r+i)}{p+1}\right\rceil.
\]
The cyclic-cover valuation formula gives
\[
 \psi_*L^{-r}
       =\bigoplus_{i=0}^{p}\OO_P(-m_i(r))\otimes N^{pi-r}.
\]
For $r\geq1$ all the relative degrees are negative.  Thus $f_*L^{-r}=0$, and the contributions to its first cohomology are
\[
 H^0\!\left(C,
   \bigl(\operatorname{Sym}^{m_i(r)-2}V\bigr)^\vee
                        \otimes N^{pi-r-(p+1)}\right)
\]
for indices with $m_i(r)\geq2$.  The dual symmetric power has a filtration with quotients $N^{-(p+1)j}$.  When $r\geq p^2-p$, even the largest additional exponent satisfies
\[
                    pi-r-(p+1)\leq p^2-r-p-1\leq-1.
\]
All contributions therefore vanish.

For the one-form bound, the exact Tango divisor identity gives
$\Omega_C^1\simeq T^p$.  The logarithmic cotangent bundle on the ruled surface has the exact sequence
\[
 0\longrightarrow\pi^*\Omega_C^1(G)=\OO_P(p)
 \longrightarrow\Omega_P^1(\log(E_0+G))
 \longrightarrow\Omega_{P/C}^1(E_0)=\OO_P(-1)\otimes T
 \longrightarrow0.
\]
Locally $G$ has equation $z^p-f_i$, with $df_i$ a nonvanishing base differential.  Allowing its logarithmic pole enlarges exactly the base differential line by $G$; its relative derivative is zero.  The disjoint infinity section enlarges exactly the relative differential line by $E_0$.  This proves the sequence directly, including in characteristic two.
The tame cyclic cover gives
$\Omega_S^1\hookrightarrow\psi^*\Omega_P^1(\log(E_0+G))$.
After tensoring with $L^{-r}$ and using the cyclic eigensheaves, the relative graded terms always have negative fibre degree.  The other terms are
\[
                    \OO_P(p-m_i(r))\otimes N^{pi-r}.
\]
At $r=p+1$, only $i=0$ has nonnegative fibre degree, and this term is $N^{-(p+1)}$, with no sections.  For $r>p+1$, every fibre degree is negative.  This proves the one-form vanishing on $S$.

For the multiroot cover, its algebra is the sum of $L^{-\sum a_i}$.  Projection formula carries the negative $H^1$ vanishing to $X$.  The logarithmic inclusion for its simple-normal-crossing branch divisor $D=\sum D_i$ bounds the twisted one-forms by the spaces
\[
 H^0\!\left(S,\Omega_S^1(\log D)\otimes L^{-(r+\sum a_i)}\right).
\]
The residue targets on each $D_i$ are negative powers of an ample line bundle, so have no sections.  Each displayed space equals the corresponding untwisted logarithmic-free one-form space, which vanishes in the asserted range.  This proves both bounds on $X$.
\end{proof}

\begin{proposition}[An integral product criterion using the relative de Rham complex]
\label{prop:OC-product-saturation-criterion}
\status{Proved; independently audited by the root and reduction-descent agents.}
Let $X$ be a smooth projective connected surface over an algebraically closed field of characteristic $p$, let $A$ be very ample, and put
\[
 Z=X\times\mathbf P^1,\qquad B=A\boxtimes\OO(m),\qquad m\geq1.
\]
For every smooth $Y\in|B^d|$, the integral map
$H^2_{\cris}(Z/W)\to H^2_{\cris}(Y/W)$ is injective.
If in addition
\begin{equation}\label{eq:OC-product-vanishing-hypotheses}
 H^1(X,A^{-jd})=0\quad(j=1,2,3),
 \qquad H^0(X,\Omega_X^1\otimes A^{-2d})=0,
\end{equation}
then its cokernel is torsion-free and it identifies the full crystalline degree-two torsion subgroups.
\end{proposition}
\begin{proof}
Let $i:Y\hookrightarrow Z$ and consider the exact sequence of complexes
\[
 0\longrightarrow K^\bullet
 \longrightarrow\Omega_Z^\bullet
 \longrightarrow i_*\Omega_Y^\bullet\longrightarrow0,
 \qquad K^j=\Omega_Z^j(\log Y)(-Y).
\]
The formula for the kernel follows in local coordinates with a smooth equation $s=0$ for $Y$: it is generated by $s\Omega_Z^j$ and $ds\wedge\Omega_Z^{j-1}$.  In particular $K^0=\OO_Z(-Y)$, and the residue sequences are
\begin{equation}\label{eq:OC-product-residue}
 0\to\Omega_Z^j(-Y)\to K^j\to\Omega_Y^{j-1}(-Y)\to0.
\end{equation}

K\"unneth gives, for every positive $r$,
\[
 H^1(Z,B^{-r})=0,\qquad
 H^2(Z,B^{-r})=H^1(X,A^{-r})\otimes H^1(\mathbf P^1,\OO(-mr)),
\]
\[
 H^0(Z,\Omega_Z^1\otimes B^{-r})=0,
 \qquad H^0(Z,\Omega_Z^2\otimes B^{-r})=0,
\]
and
\begin{equation}\label{eq:OC-product-omega-H1}
 H^1(Z,\Omega_Z^1\otimes B^{-r})
  =H^0(X,\Omega_X^1\otimes A^{-r})
                 \otimes H^1(\mathbf P^1,\OO(-mr)).
\end{equation}
For the last identity, the summand from $\Omega^1_{\mathbf P^1}$ contributes zero because the needed $H^0$ of $A^{-r}$ or of $\Omega^1_{\mathbf P^1}(-mr)$ vanishes.  These are direct tensor-product cohomology calculations.

Since $\OO_Y(-Y)$ is antiample, it has no global sections.  The $j=1$ sequence in \eqref{eq:OC-product-residue} gives $H^0(K^1)=0$.  Along with $H^1(K^0)=H^1(Z,B^{-d})=0$, this kills both total-degree-one terms of the Hodge hypercohomology spectral sequence of $K^\bullet$.  Therefore
\[
                         \mathbb H^1(Z,K^\bullet)=0
\]
for every $d$.  The de Rham long exact sequence shows that
$H^1_{\dR}(Z)\to H^1_{\dR}(Y)$ is injective.

The first crystalline groups are free of the same rank by rational weak Lefschetz.  Their reductions modulo $p$ embed in first de Rham cohomology, naturally for restriction; hence their map modulo $p$ is injective.  A map of finite free modules over $W$ of the same rank which is injective modulo $p$ is an isomorphism.  In the first universal-coefficient diagram the left map is consequently an isomorphism, and the snake lemma identifies
\[
 \ker\{H^2_{\cris}(Z/W)[p]\to H^2_{\cris}(Y/W)[p]\}
    =\ker\{H^1_{\dR}(Z)\to H^1_{\dR}(Y)\}=0.
\]
Rational weak Lefschetz makes the full second crystalline kernel torsion.  A nonzero finite $p$-primary torsion module has a nonzero element killed by $p$, so that full kernel is zero.  This proves the first assertion for all $d$.

Assume now \eqref{eq:OC-product-vanishing-hypotheses}.  The preceding identities and divisor sequences give
\[
 H^2(K^0)=0,\qquad
 H^1(Y,\OO_Y(-Y))=H^1(Y,\OO_Y(-2Y))=0,
\]
\[
                     H^1(Z,\Omega_Z^1(-2Y))=0.
\]
For example the first line for $Y$ follows from
$0\to\OO_Z(-2Y)\to\OO_Z(-Y)\to\OO_Y(-Y)\to0$
and its version at $-2Y$, using respectively the negative $H^2$ groups at $2d$ and $3d$.
The restriction sequence for $\Omega_Z^1(-Y)$ then gives
$H^0(Y,\Omega_Z^1|_Y(-Y))=0$.
The twisted conormal sequence
\[
 0\to\OO_Y(-2Y)\to\Omega_Z^1|_Y(-Y)\to\Omega_Y^1(-Y)\to0
\]
therefore gives $H^0(Y,\Omega_Y^1(-Y))=0$.
Together with $H^0(Z,\Omega_Z^2(-Y))=0$, the $j=2$ residue sequence proves $H^0(K^2)=0$.

Thus the only possible total-degree-two term on the first Hodge page is $H^1(K^1)$.  It is canonically $H^1(Z,\Omega_Z^1(-Y))$, since both $H^0(\OO_Y(-Y))$ and $H^1(\OO_Y(-Y))$ vanish.  Its outgoing differential is injective.  To see this without assuming the vanishing of that group, compose it with the residue map:
\[
 H^1(Z,\Omega_Z^1(-Y))\simeq H^1(K^1)
 \xrightarrow{d_1}H^1(K^2)
 \longrightarrow H^1(Y,\Omega_Y^1(-Y)).
\]
This composite, up to the harmless residue sign, is the ordinary restriction map.  Locally its identity is exactly
\[
                      d(s\eta)=ds\wedge\eta+s\,d\eta;
\]
the residue of the first term is $\eta|_Y$ with its conormal twist.
The restriction is injective: its factor through
$H^1(Y,\Omega_Z^1|_Y(-Y))$ is injective because
$H^1(Z,\Omega_Z^1(-2Y))=0$, and the next map is injective because
$H^1(Y,\OO_Y(-2Y))=0$ in the twisted conormal sequence.
Consequently $d_1$ is injective.  All total-degree-two groups on the second Hodge page are zero, and hence
\[
                          \mathbb H^2(Z,K^\bullet)=0.
\]
The de Rham long exact sequence now gives an isomorphism on $H^1_{\dR}$ and an injection on $H^2_{\dR}$.

We finish the integral argument with the torsion terms retained.  Write
$M=H^2_{\cris}(Z/W)$ and $N=H^2_{\cris}(Y/W)$.  We already proved $M\hookrightarrow N$.  The first universal-coefficient diagram, with its isomorphisms on the first crystalline lattice and on first de Rham cohomology, gives
\[
                              M[p]\xrightarrow{\sim}N[p].
\]
The second universal-coefficient diagram and the de Rham injection give
$M/pM\hookrightarrow N/pN$.
For $C=N/M$, the exact Tor sequence includes
\[
 M[p]\longrightarrow N[p]\longrightarrow C[p]
                 \longrightarrow M/pM\longrightarrow N/pN.
\]
The first map is surjective and the last map injective, so $C[p]=0$.
Thus $C$ is torsion-free.  The exact sequence $0\to M\to N\to C\to0$ then identifies the full torsion subgroups of $M$ and $N$.  This uses both universal-coefficient diagrams; injection modulo $p$ alone would not suffice for this conclusion.

The first-crystalline freeness and natural universal-coefficient sequence are as in
\cite[II, Proposition~3.11 and Remark~3.11.2]{Illusie1979DRW}; rational weak Lefschetz is
\cite[Appendix~9, Corollary~9.3, pp.~63--64, version~3]{NakkajimaYobuko2022}.
\end{proof}

\begin{theorem}[Exact degrees for the characteristic-two integral product counterexample]
\label{thm:OC-product-exact-degrees}
\status{Proved; independently audited by the root and reduction-descent agents.}
Take the characteristic-two surface $(X,A)$ of
Corollary~\ref{cor:OC-surface-exact-degrees}, and the fixed threefold embedding
\[
 m=591\cdot9^{556}+2,\qquad Z=X\times\mathbf P^1,
 \qquad B=A\boxtimes\OO(m).
\]
For every smooth degree-$d$ hypersurface section $Y_d$ in this embedding:
\begin{enumerate}
\item The map $H^2_{\cris}(Z/W)\to H^2_{\cris}(Y_d/W)$ is injective for every $d\geq1$.
\item Its cokernel is torsion-free if and only if $d\geq2$.
\item Its map on the full crystalline degree-two torsion subgroups is an isomorphism if and only if $d\geq2$.
\end{enumerate}
Thus the exact degree range requested in A.1(a) and A.1(b) is determined for this counterexample.
\end{theorem}
\begin{proof}
The all-degree injectivity follows from
Proposition~\ref{prop:OC-product-saturation-criterion} without extra hypotheses.
For every $d\geq2$, the already computed surface cohomology gives
$H^1(X,A^{-jd})=0$ for $j=1,2,3$.
Also $2d\geq4\geq3$, so
Lemma~\ref{lem:OC-surface-uniform-negative-bounds} gives
$H^0(X,\Omega_X^1\otimes A^{-2d})=0$.
The criterion proves both positive assertions for every such $d$.
At $d=1$, Theorem~\ref{thm:OC-explicit-integral-counterexample} proves nonzero cokernel torsion and failure of torsion surjectivity for every smooth hyperplane.  This supplies both converses.
\end{proof}

\begin{corollary}[A uniform degree bound for all-prime product examples]
\label{cor:OC-all-prime-product-degrees}
\status{Proved; independently audited by the root and reduction-descent agents.}
For every prime $p$, the threefold of
Corollary~\ref{cor:OC-all-prime-integral-counterexample} has injective integral second-crystalline restriction in every positive degree.  Its cokernel is torsion-free and its full torsion groups are identified for every
\[
                              d\geq p^2-p.
\]
Both latter assertions fail in degree one.  At $p=2$ this bound is exact by
Theorem~\ref{thm:OC-product-exact-degrees}.  No claim about the intervening degrees $2\leq d<p^2-p$ is made for the additional examples at odd primes.
\end{corollary}
\begin{proof}
For $d\geq p^2-p$, the first surface bound supplies all three negative $H^1$ vanishings.  Also $2d\geq p+1$ for every prime, so the second surface bound supplies the one-form vanishing in
\eqref{eq:OC-product-vanishing-hypotheses}.  Apply the product criterion and the existing degree-one counterexample.
\end{proof}


\begin{remark}[Exact arithmetic and torsion checks]
\label{rem:OC-product-degree-checks}
The certificate \texttt{A\_product\_degree\_checks.py} checks the boundary eigensheaf inequalities for eleven primes, the characteristic-two threshold, and $113$ homomorphisms of small cyclic primary modules.  It also records explicit counterexamples to each insufficient shortcut in the final Tor argument.  The recorded run returned PASS in $0.0006374399999913294$ seconds.  The all-prime inequalities and the geometric conclusions are proved above; finite enumeration is not used to extrapolate them.
\end{remark}
